\documentclass[../DoAn.tex]{subfiles}
\begin{document}

\section{Kết luận}
\label{section:6.1}
Sau quá trình nghiên cứu lý thuyết, thiết kế hệ thống và tiến hành thực nghiệm lập trình chi tiết, đồ án đã hoàn thành đầy đủ các mục tiêu và nhiệm vụ đặt ra trong phiếu giao đề tài tốt nghiệp. Các kết quả nổi bật đạt được bao gồm:
\begin{enumerate}
    \item \textbf{Về mặt mô hình AI:} Thiết kế và huấn luyện thành công mô hình học sâu lai (Hybrid) kết hợp tối ưu giữa chỉ số hình học khuôn mặt từ MediaPipe FaceMesh (tính toán EAR, MAR, Pitch để phát hiện ngủ gật, ngáp và mất tập trung) và mô hình CNN MobileNetV3-Large phân loại cảm xúc tiêu cực của tài xế. Mô hình CNN đạt độ chính xác phân loại \textbf{88.4\%} trên tập dữ liệu kiểm thử thực tế.
    \item \textbf{Về mặt hệ thống:} Xây dựng hoàn chỉnh giải pháp từ biên lên trung tâm (Edge-to-Cloud). Phía Client trên xe xử lý video đa luồng mượt mà ở tốc độ 30 FPS, phát cảnh báo âm thanh cục bộ bất đồng bộ nhanh chóng dưới 80ms. Phía Server điều hành trung tâm sử dụng Flask Web Server tiếp nhận dữ liệu ổn định, lưu trữ SQLite an toàn với cơ chế lọc trùng debounce 10 giây, và đồng bộ dữ liệu Dashboard thời gian thực thông qua cơ chế Server-Sent Events (SSE).
    \item \textbf{Về thuật toán tối ưu:} Đề xuất và cài đặt thành công bộ lọc làm mượt xác suất cảm xúc ProbSmoother kết hợp bộ tích lũy cảm xúc MultiEmotionScorer giúp triệt tiêu hiện tượng flickering nhãn cảnh báo; thuật toán tự động hiệu chỉnh ngưỡng EAR động theo sinh học khuôn mặt tài xế; và thuật toán tự động chấm điểm an toàn Safety Score phục vụ đánh giá xếp hạng tài xế chuyên nghiệp.
\end{enumerate}

\noindent \textbf{Hạn chế của hệ thống:}
\begin{itemize}
    \item Hệ thống sử dụng camera ánh sáng thường (RGB) nên hoạt động kém hiệu quả trong điều kiện cabin hoàn toàn tối vào ban đêm nếu không có đèn hồng ngoại phụ trợ.
    \item Hệ thống chưa nhận diện được các hành vi vi phạm khác như tài xế hút thuốc lá, nghe điện thoại bằng tay hoặc không thắt dây an toàn.
\end{itemize}

\section{Hướng phát triển}
\label{section:6.2}
Để hoàn thiện hệ thống hướng tới thương mại hóa thực tế, các hướng nghiên cứu phát triển tiếp theo bao gồm:
\begin{enumerate}
    \item \textbf{Nâng cấp phần cứng camera:} Tích hợp camera hồng ngoại (IR Camera) kết hợp đèn LED hồng ngoại bước sóng 850nm để hệ thống có khả năng giám sát tài xế ổn định trong mọi điều kiện ánh sáng, kể cả ban đêm hoàn toàn.
    \item \textbf{Tối ưu hóa mô hình nhúng:} Thực hiện chuyển đổi mô hình PyTorch sang định dạng ONNX hoặc TensorRT, áp dụng các kỹ thuật lượng tử hóa trọng số (Quantization từ FP32 xuống INT8) để giảm thiểu kích thước mô hình và tăng tốc độ suy luận, giúp triển khai mượt mà trên các dòng chip nhúng siêu rẻ như Raspberry Pi 4/5 hoặc Jetson Nano.
    \item \textbf{Mở rộng tính năng giám sát hành vi:} Phát triển thêm các module AI phát hiện vật thể (như YOLOv8-Pose hoặc YOLOv8-Object) để phát hiện và cảnh báo các hành vi vi phạm khác của tài xế bao gồm: nghe điện thoại bằng tay, hút thuốc lá trong cabin, không thắt dây an toàn hoặc không đặt tay lên vô lăng.
    \item \textbf{Tích hợp bản đồ và cảnh báo từ xa:} Phát triển thêm tính năng bản đồ số GPS trên Dashboard quản trị trung tâm, hiển thị vị trí thời gian thực của các xe vi phạm lỗi nặng để điều phối viên có thể liên lạc cứu hộ hoặc cảnh báo trực tiếp qua đàm thoại IP kịp thời.
\end{enumerate}

\end{document}
